\documentclass[language=english,debug=true,a4paper,twoside,12pt]{lmc}
% Options: same options as standard book class
%and two new keyval options:
% - boolean debug, that shows/hide pageframe, line number and
% custom pages with debug info at the end of the document  (Try it!)
% - language, to apply to babel package. We support correctly 
% english and portuguese

% Compilation:
% I recomend using Miktex as the compiler and VSCode + Latex Workshop 
% extension and the IDE. To compile this template within the suggested
% framework, just go to the command, palette (ctrl+shift+p), search for Latex Workshop: Build with recipe, and use the recipe PDFLatex + Biber + PDFLatex*2
% If you are using other IDEs, you will have to search how to use Biber there, 
% since ABNT-style citation needs it instead of Bibtex.
% I haven't tried it in Overleaf yet.

% Useful tips:
%   Using subfile package, you can compile either individual subfiles or the 
% whole main document.
%   This template is capable of bidirectional navigation (see synctex=1 in
% class definition).
%   In vscode, ctrl+click on PDF to go to source code and select text in source
% code and press ctrl+alt+j to go to PDF.
%   Note that the subfiles will map to the PDF created when compiling the 
% subfile only.

% REMARK: For figures and tables add label after caption. Note that if caption is positioned above the figure/table, the caption is positioned above

% If you want to automatically use ABNT citations, you MUST use biblatex with biber backend. You can change those configurations as you wish.
\usepackage[style=abnt-numeric,backend=biber,citecounter=true,backrefstyle=three]{biblatex}
\usepackage{csquotes}
%Adds a period and boldface the reference number ONLY WHEN IN REFERENCE LIST, for style=abnt-numeric
\DeclareFieldFormat{labelnumber}{\ifbibliography{\mkbibbold{#1\addperiod}}{#1}}

% Adds the .bib sources
\addbibresource{\subfix{custom_bibliography.bib}}

% Use subfiles + xr-hyper (loaded by class) to get correct correspondences!
% I recomend adding M- to all \refs, \crefs, etc, to guarantee correct hyperrefs. Also, you will need to compile the whole main at least onece before making it work correctly
\externaldocument[M-]{\subfix{main}}

\input{\subfix{macros.tex}}

% In this template, all the commands that are introduced
% in this custom class are used for the cover and titlepage.
% PLEASE FILL ALL THE METADATA BELOW
\author{My Name}
\title{My book}
\City{My city}
\date{\the\year}
\Docversion{Original Version}
\Advisor{My advisor}  
\ConcentrationArea{My concentration area}
\School{Polytechnic School of the University of São Paulo}
\Department{Department of Structural and Geotechnical Engineering}
\TitlePageText{Thesis submitted to the \theSchool{}, to obtain 
                the degree of Doctor of Science in the Postgraduate Program in Civil Engineering.\par
                Concentration area:\par
                \textit{\theConcentrationArea{}}\par
                Advisor: \par \textit{\theAdvisor}}

\begin{document}
%    
%%============= pre-textual content =================

%modify locally the plain pagestyle, so that it
%gets used in the TOC, TOL and TOF correctly
%========= COVER AND TITLE PAGE ===================
% CAPA
\cover 
% FOLHA DE ROSTO + ficha catalográfica, if available
\TitlePage{}
% OR
%\TitlePage{ficha_catalografica.pdf}

%===== EVALUATION SHEET ========
%\includepdf{my_evaluation_sheet.pdf} evaluation sheet


%========= DEDICATION ===================
\renewcommand{\abstractname}{Dedication}
\begin{abstract}
    You can \verb!\renewcommand! in \verb!\abstractname! to create unnumbered environements for the pretextual elements!
\end{abstract}

%========= ACKNOWLEDGEMENTS ===================
\renewcommand{\abstractname}{Acknowledgements}
\begin{abstract}
    Acknowledgements ....
\end{abstract}

%================ EPIGRAPH ====================
\cleardoublepage

\vspace*{\fill}
\epigraph{
    "(...) \textit{some fancy text} (...)"
    \vspace{6pt}\\
    "(...) \textit{algum texto legal} (...)" 
    \vspace{6pt}
}{--- \textit{An important guy}, somehwere, someday.}
\vspace*{\fill}

%========= RESUMO [PT-BR] ===================
\renewcommand{\abstractname}{Resumo}
\begin{abstract}
    
    %\include{abstract.tex}
    Insira resumo aqui...
    %
    \\[3\baselineskip]
    %
    \textbf{Palavras-chave} -- Word, Word, Word, Word, Word.
\end{abstract}

%========= Abstract [English] ===================

\renewcommand{\abstractname}{Abstract}
\begin{abstract}
    %\include{abstract.tex}
    Insert abstract here...
    %
    \\[3\baselineskip]
    %
    \textbf{Keywords} -- Word, Word, Word, Word, Word. 
\end{abstract}

% === Tables of contents, figures and tables ===
\tableofcontents
\listoffigures
\listoftables

% ========== BEGIN OF TEXTUAL CONTENT ===================
% You can add subfiles with the actual content!
\subfile{chapters/chapter_one/chapter1.tex}
\subfile{chapters/chapter_two/chapter2.tex}
\subfile{chapters/chapter_three/chapter3.tex}

% Or just do regular latex commands!
\chapter{My fourth chapter}
\lipsum[9-9]
\chapter{My fifth chapter}
\lipsum[9]
\section{First section of fifth chapter}
\subsection{First subsection of first section of fifth chapter}
\subsubsection{First subsubsection of first subsection of first section of fifth chapter}
\lipsum[15]
\section{Second section of second chapter}
\lipsum[15]
See \cite{timoshenko_theory_1985}
\begin{equation}
    a = b + c
    \label{eq:sum}
\end{equation}
\lipsum[3-20]
\chapter{My sixth chapter}
\lipsum[9]


%===============================BIBLOGRAPHY===============================

\cleardoublepage %ends the current page and causes all figures and tables that have so far appeared in the input to be printed. 
\printbibliography

%=================================APPENDIX=================================
\appendix
\chapter{First appendix}
\section{Section inside appendix}
\subsection{Sub Section inside appendix}
\subsubsection{Sub Sub Section inside appendix}
\chapter{Second appendix with a very a very a very a very a very a very a very a very long title}

\lipsum[9]

%==================================ANNEX===================================
\annex
\chapter{First annex}
\section{Section inside annex}
\subsection{Sub Section inside annex}
\subsubsection{Sub Sub Section inside annex}
\begin{equation}
    a = b^2 + c^2
\end{equation}
\chapter{Second annex}


\DEBUG{
\chapter{This is a DEBUG chapter}
This is only shown if debug = true
}


\end{document}
